\documentclass{article}

\PassOptionsToPackage{numbers}{natbib}
\usepackage[preprint]{neurips_2026}

\usepackage[utf8]{inputenc}
\usepackage[T1]{fontenc}
\usepackage{hyperref}
\usepackage{url}
\usepackage{booktabs}
\usepackage{amsfonts}
\usepackage{amsmath}
\usepackage{amssymb}
\usepackage{nicefrac}
\usepackage{microtype}
\usepackage{xcolor}
\usepackage{graphicx}
\usepackage{enumitem}
\usepackage{listings}

\lstset{
  basicstyle=\small\ttfamily,
  breaklines=true,
  frame=single,
  xleftmargin=1em,
}

\title{Lazy Row-Group Loading in \texttt{ChanghunDataset}:\\
  Design, Trade-offs, and the LVN Bottleneck}

\author{Changhun Kim}

\begin{document}
\maketitle

%==============================================================================
\section{Why Lazy Loading Is Needed}
%==============================================================================

In the eager (default) loading mode every row of the Parquet file is read,
decoded, and converted to PyTorch tensors at dataset initialisation time.
For small cases this is fine, but for large graphs the memory requirement is
prohibitive.

A single row of an $N$-bus dataset contains:
\begin{itemize}[noitemsep]
  \item \texttt{Y\_matrix}: $N \times N$ complex64 $\;\to\;$ $8N^2$ bytes.
  \item \texttt{Y\_Lines}: $\binom{N}{2}$ complex64 entries.
  \item \texttt{u\_start}, \texttt{u\_newton}, \texttt{S\_start},
    \texttt{S\_newton}: $N$-vectors.
\end{itemize}

For \textbf{case1354} ($N=1354$), one row's \texttt{Y\_matrix} alone is
$8 \times 1354^2 \approx 14.7\,\mathrm{MB}$.  Storing 36{,}000 rows eagerly
would require $\sim\!530\,\mathrm{GB}$ of RAM — completely infeasible.
Lazy row-group mode was introduced to solve this.

%==============================================================================
\section{Parquet Row Groups: The Building Block}
%==============================================================================

A Parquet file is physically organised into \emph{row groups} — independent
chunks that can be read from disk without touching the rest of the file:

\begin{lstlisting}
LVN_36k.parquet  (1.1 GB on disk)
  Row Group 0   :  20 rows  (~620 KB)
  Row Group 1   :  20 rows
  ...
  Row Group 1799:  20 rows
\end{lstlisting}

Because each row group is self-contained, the dataset can read only the
row groups it currently needs and discard the rest.

%==============================================================================
\section{Eager Mode: Simple, Fast, Memory-Heavy}
%==============================================================================

\begin{lstlisting}[language=Python]
ChanghunDataset(path, lazy_row_groups=False)   # default
\end{lstlisting}

\paragraph{Initialisation.}
The entire Parquet file is read once; every row is decoded and converted to
tensors.  The result is stored as \texttt{self.rows = [dict\_0, dict\_1,
\ldots]}, fully resident in RAM.

\paragraph{Item access.}
\texttt{\_\_getitem\_\_(idx)} is a single list index operation --- microsecond
latency, the GPU is never starved.

\paragraph{Trade-off.}
RAM usage equals the full dataset size.  Unusable for large-$N$ cases.

%==============================================================================
\section{Lazy Row-Group Mode: Chunked Read with LRU Cache}
%==============================================================================

\begin{lstlisting}[language=Python]
ChanghunDataset(path, lazy_row_groups=True, row_group_cache_size=2)
\end{lstlisting}

\subsection{Initialisation (\texttt{\_build\_lazy\_row\_group\_index})}

Only the Parquet \emph{metadata} is scanned --- no payload data is decoded.
For each row group the initialiser records:
\begin{itemize}[noitemsep]
  \item its row-group index and the number of rows it contains;
  \item which of those rows passed the divergence filter (a lightweight
    read of \texttt{u\_newton} only).
\end{itemize}

The result is a compact index table \texttt{self.\_row\_group\_metas}
(tens of MB at most) and an empty LRU cache
\texttt{self.\_row\_group\_cache = OrderedDict()}.

\subsection{Item Access at Training Time}

\texttt{\_\_getitem\_\_(idx)} proceeds in three steps:

\begin{enumerate}[noitemsep]
  \item \textbf{Map index to row group.}
    Binary search over \texttt{\_row\_group\_metas} finds the row group $g$
    that contains row \texttt{idx}, and the local offset $r$ within $g$.
  \item \textbf{Cache lookup.}
    \begin{itemize}[noitemsep]
      \item \textit{Cache hit}: row group $g$ is already in
        \texttt{\_row\_group\_cache} $\to$ return row $r$ immediately.
      \item \textit{Cache miss}: read row group $g$ from disk, decode all
        binary columns, convert to tensors, insert into cache.
        If the cache exceeds \texttt{row\_group\_cache\_size}, evict the
        least-recently-used entry.
    \end{itemize}
  \item \textbf{Return} row $r$ from the (now-populated) cache entry for
    group $g$.
\end{enumerate}

\subsection{Cache Visualisation}

With \texttt{row\_group\_cache\_size=2} the in-memory footprint at any
moment is at most two row groups (e.g.\ 40 rows for a 20-row/group file):

\begin{lstlisting}
  _row_group_cache (OrderedDict, max 2 entries)
  +--------------------+  +--------------------+
  |  Row Group 617     |  |  Row Group 1203    |
  |  20 decoded rows   |  |  20 decoded rows   |
  +--------------------+  +--------------------+
         ^-- most recent

  Next batch requests Row Group 99:
    -> evict Row Group 1203 (oldest)
    -> load Row Group 99 from disk
\end{lstlisting}

%==============================================================================
\section{Comparison: Eager vs.\ Lazy}
%==============================================================================

\begin{table}[h]
\centering
\caption{Eager vs.\ lazy row-group mode.}
\label{tab:eager_vs_lazy}
\begin{tabular}{lll}
\toprule
Aspect & Eager & Lazy row-group \\
\midrule
Init time        & Long (decode everything)   & Short (metadata only) \\
Init RAM         & Full dataset               & Index only ($\sim$tens MB) \\
\texttt{\_\_getitem\_\_} latency
                 & Microseconds               & Disk I/O + decode on miss \\
Sequential access & Fast                      & Fast (many cache hits) \\
Random / shuffled & Fast                      & Slow (many cache misses) \\
RAM at runtime   & Full dataset               & $\propto$ cache size \\
Best for         & Small $N$, small datasets  & Large $N$ ($\geq$1000 buses) \\
\bottomrule
\end{tabular}
\end{table}

%==============================================================================
\section{Cache Hit Rate Analysis}
%==============================================================================

Under random shuffling, the expected cache hit rate is approximately:

\begin{equation}
  \text{hit rate}
  \;\approx\;
  \frac{\text{cache size (row groups)}}{\text{total row groups}}
  \label{eq:hitrate}
\end{equation}

Table~\ref{tab:hitrate} shows how this plays out for different cases with
the same 36{,}000-row dataset and \texttt{cache=2}.

\begin{table}[h]
\centering
\caption{Approximate cache hit rate under random shuffling
  (\texttt{cache=2}, 36k rows).}
\label{tab:hitrate}
\begin{tabular}{lrrl}
\toprule
Case & Row groups & Hit rate & Consequence \\
\midrule
case300 (36k)  & 6    & $2/6    = 33\%$ & Tolerable \\
case1354 (36k) & 120  & $2/120  = 1.6\%$ & Slow \\
LVN (36k)      & 1800 & $2/1800 = 0.1\%$ & \textbf{Catastrophic} \\
\bottomrule
\end{tabular}
\end{table}

For the LVN 36k dataset, almost every \texttt{\_\_getitem\_\_} call triggers
a full disk read and decode cycle.  This is the root cause of the
$\sim\!15$--$16\,\mathrm{min/epoch}$ wall time observed in the LVN training
runs.

The underlying reason is that the LVN Parquet was written with a very small
row-group size (20 rows/group), producing 1{,}800 groups for 36{,}000 rows.
By contrast, case300 uses 6{,}000 rows/group, so the same 36{,}000 rows fit
in only 6 groups and the cache works as intended.

%==============================================================================
\section{Remedies}
%==============================================================================

\subsection{Option A: Increase Cache Size (Immediate)}

\begin{lstlisting}
--row_group_cache_size 200
\end{lstlisting}

Raising the cache from 2 to 200 row groups raises the hit rate from
$0.1\%$ to $200/1800 \approx 11\%$ and reduces disk I/O by roughly
$10\times$.  The RAM cost is $200 \times 20\,\text{rows} \times
(\text{row size})$; for LVN with $N \approx 722$ buses this is
approximately $200 \times 20 \times 4\,\text{MB} \approx 16\,\text{GB}$
of additional RAM.

\subsection{Option B: Re-write the Parquet with Larger Row Groups
  (Root Fix)}

Re-generating the LVN Parquet in \texttt{ScenarioSynthesis\_PPC} with a
larger \texttt{row\_group\_size} (e.g.\ 2{,}000):

\begin{lstlisting}[language=Python]
table.write_to_dataset(..., row_group_size=2000)
\end{lstlisting}

reduces the number of row groups from 1{,}800 to 18, so a cache of size 2
already covers $11\%$ of the dataset --- the same hit rate as Option~A but
without the extra RAM.  This is the preferred long-term fix.

\subsection{Option C: Sparse In-Memory Format}

Store the dataset in RAM using a sparse tensor representation so that even
case1354-scale grids fit in a reasonable memory budget.  Lazy mode then
only needs to perform the tensor type conversion, not disk I/O.

%==============================================================================
\section{Intuitive Analogy}
%==============================================================================

\textbf{Eager mode} is like spreading an entire book flat on a desk: any
page is one finger-movement away, but the desk must be large enough to hold
the whole book.

\textbf{Lazy mode} is like keeping the book on a shelf and bringing only a
few chapters to the desk at a time.  This works well when:
\begin{itemize}[noitemsep]
  \item chapters are large (few of them $\to$ high cache hit rate), and
  \item access is somewhat sequential.
\end{itemize}
It fails when chapters are very thin (many of them) and reading order is
random --- exactly the situation with LVN's 1{,}800-group Parquet under
shuffled training.

%==============================================================================
\section{Summary}
%==============================================================================

\begin{itemize}[noitemsep]
  \item Lazy row-group mode avoids RAM OOM for large-$N$ graphs by reading
    only the Parquet row groups that are currently needed, keeping the rest
    on disk behind an LRU cache.
  \item Cache hit rate $\approx$ cache size / total row groups under random
    access.  A small cache is only effective when row groups are large.
  \item The LVN 36k Parquet has 1{,}800 row groups of 20 rows each.  With
    cache size 2, the hit rate is $0.1\%$, making nearly every training step
    a disk read.
  \item The root cause is the small \texttt{row\_group\_size} used when the
    LVN Parquet was written.  The fix is either to increase the cache size
    at training time or to re-write the Parquet with larger row groups.
\end{itemize}

\end{document}
